\documentclass[12pt]{article}

\setlength{\parindent}{0pt}
\setcounter{secnumdepth}{3}
\usepackage[margin=1in]{geometry}

% packages
\usepackage{amsmath, amsfonts, amsthm, amssymb}
\usepackage{graphicx, float}
\usepackage{mathtools}
\usepackage{titlesec}
\usepackage{interval}
\usepackage{hyperref}
\usepackage{siunitx}
\usepackage{titling}
\usepackage{vwcol}
\usepackage{setspace}
\usepackage{empheq}
\usepackage{cancel}
\usepackage{esdiff}
\usepackage{multicol}
\usepackage{mdframed}
\usepackage{esdiff}
\usepackage{tikzsymbols}
\usepackage{multicol}
\usepackage{tikz}
\usepackage{varwidth}
\usepackage{pgfplots}
\usepackage{fancyhdr}
\usepackage{enumitem}
\usepackage{tcolorbox}

\intervalconfig {
	soft open fences
}

% header
\pagestyle{fancy}
\fancyhf{}
\lhead{Ethan Fan}
\rhead{AP Calculus BC}
\cfoot{\thepage}

% section format
\titleformat{\section}
  {\large \bfseries}
  {}
  {0em}
  {}
\titlespacing*{\section}{0pt}{0.5em}{0.3em}

\titleformat{\subsection}[runin]
  {\bfseries}
  {}
  {0em}
  {}
\titlespacing*{\subsection}{0pt}{0pt}{0em}

\titleformat{\subsubsection}[runin]
  {\itshape}
  {(\roman{subsubsection})}
  {0em}
  {}

% commands
\newcommand{\alignedintertext}[1]{%
  \noalign{%
    \vskip\belowdisplayshortskip
    \vtop{\hsize=\linewidth#1\par
    \expandafter}%
    \expandafter\prevdepth\the\prevdepth
  }%
}

\newcommand*{\paren}[1]{\ensuremath\left(#1\right)}
\newcommand*{\limit}[2][x]{\ensuremath{\displaystyle\lim_{#1\to#2}}}
\newcommand*{\Limit}[3][x]{\ensuremath{\displaystyle\lim_{#1\to#2}\left[#3\right]}}
\newcommand*{\deriv}[1][x]{\ensuremath{\dfrac{\mathrm{d}}{\mathrm{d}#1}}}
\newcommand*{\Deriv}[2][x]{\ensuremath{\dfrac{\mathrm{d}}{\mathrm{d}#1}\left[#2\right]}}
\newcommand{\ddx}{\dfrac{d}{dx}}
\newcommand{\dydx}{\dfrac{dy}{dx}}
\newcommand{\dydt}{\dfrac{dy}{dt}}
\newcommand{\dxdt}{\dfrac{dx}{dt}}
\newcommand*{\iinteg}[2][x]{\ensuremath{\displaystyle\int #2\;\mathrm{d}#1}}
\newcommand*{\dinteg}[4][x]{\ensuremath{\displaystyle\int_{#2}^{#3}#4\;\mathrm{d}#1}}
\newcommand*{\abs}[1]{\ensuremath{\left|#1\right|}}
\newcommand*{\eps}{\varepsilon}
\newcommand*{\real}{\ensuremath\mathbb{R}}
\newcommand*{\integer}{\ensuremath\mathbb{Z}}
\newcommand*{\posint}{\ensuremath\mathbb{Z^+}}
\newcommand*{\cbrt}[1]{\ensuremath{\sqrt[3]{#1}}}
\newcommand*{\lheqzero}{\ensuremath{\underset{\text{L'H}}{\overset{\left[\frac00\right]}{=}}}}
\newcommand*{\lheqinfty}{\ensuremath{\underset{\text{L'H}}{\overset{\left[\frac{\infty}{\infty}\right]}{=}}}}
\newcommand{\tor}{\text{ or }}
\newcommand{\floor}[1]{\lfloor #1 \rfloor}
\newcommand{\at}{\Big|}

\newcommand{\vem}{\vspace{0.5em}}
\newcommand{\example}[1]{\section{Example #1}}
\newcommand{\problem}[1]{\section{Problem #1}}
\newcommand{\subproblem}[1]{\subsection{(#1)} \,}

\DeclareMathOperator{\dne}{DNE}
\DeclareMathOperator{\sgn}{sgn}

\newcommand{\definition}[1]{\begin{tcolorbox}[colback=red!5!white,colframe=red!75!black,parbox=false] #1 \end{tcolorbox}}

\newcommand{\theorem}[2]{\begin{tcolorbox}[title={#1},colback=blue!5!white,colframe=blue!75!black,parbox=false] #2 \end{tcolorbox}}

\allowdisplaybreaks
% \postdisplaypenalty=100000

% document
\begin{document}

\vspace*{0cm}

% opening
\begin{center}
    {\Large \bfseries Problem Set 26} \\[0.5cm]
    {\large Rotation at Work: Volumes, Lengths, and Surfaces} \\[0.35cm]
    {\normalsize November 10, 2026}
\end{center}

% problems

\problem{1}

\subproblem{a}
\begin{gather*}
    F=\int_0^R 2\pi r \, v(r)\, dr\\
    = 2\pi k \int_0^R rR^2-r^3\, dr\\
    = 2\pi (\frac{r^2}{2}R^2-\frac{r^4}{4})\at_0^R\\
    = \boxed{\frac{\pi k}{2}R^4}
\end{gather*}

\subproblem{b}
\begin{gather*}
    F/\pi R^2 = \frac{1}{2}kR^2\\
    v(0)=kR^2 
\end{gather*}

\subproblem{c}
\begin{gather*}
    1-0.9^4=0.3439 \implies 34.49 \%\\
    \frac{1}{0.9^4} \approx 1.524 \implies 52.4\%
\end{gather*}

\problem{2}

\subproblem{c}
\begin{gather*}
    S = 2\pi b^2 + 2\pi ab \frac{\arcsin e}{e}\\
    \lim_{e \to 0} S = 2\pi b^2 + 2\pi ab  \\
    = 2\pi b(a+b)\\
    =\boxed{4\pi a^2}
\end{gather*}
The surface area approaches that of a sphere.\vem

\subproblem{d}
\begin{gather*}
    V = \frac{4}{3}\pi a b^2 = \frac{4}{3}\pi \cdot 2.85\cdot 2.2^2 \approx 57.78 \text{ cm}^3\\
    S = 2\pi \cdot 2.2^2 + 2\pi \cdot 2.85\cdot2.2 \cdot \frac{\arcsin e}{e} \approx 73.09 \text{ cm}^2\\
    V \approx 73.09  \cdot 0.035 \approx 2.56 \text{ cm}^3
\end{gather*}

\subproblem{e}
\begin{gather*}
    \frac{4}{3}\pi r^3 = 57.78 \implies r = \left(\frac{3 \cdot 57.78}{4\pi}\right)^{\frac{1}{3}}\\
    S = 4\pi r^2 \approx 72.28 \text{ cm}^2\\
    \frac{73.09 - 72.28}{72.28} \approx 1.12\%
\end{gather*}

\problem{3}

\subproblem{a}
\begin{gather*}
    d V=\frac{4}{3}\pi  \paren{ (r+dr)^3 -r^3} \\
    =\frac{4}{3}\pi (3r^2\, dr+3r\, d^2r + d^3r)\\
    = \boxed{4\pi r^2}
\end{gather*}

\subproblem{b}
\begin{gather*}
    \frac{1}{\pi a^3} e^{-2r/a} \, dv = \frac{4}{a^3}r^2 e^{-2r/a} \, dr\\
    P= \int_0^{ka} \frac{4}{a^3}r^2 e^{-2r/a} \, dr\\
    = \begin{bmatrix}
        u=r^2 & dv = e^{-2r/a} \, dr\\
        du=2r\, dr & v= -\frac{a}{2}e^{-2r/a}
    \end{bmatrix}\\
    - \frac{4}{a^3} \cdot \frac{a}{2}\cdot r^2 e^{-2r/a} \at^{ka}_0 + \frac{4}{a^2} \int_0^{ka}  r e^{-2r/a}\, dr\\
    \begin{bmatrix}
        u=r & dv = e^{-2r/a} \, dr\\
        du= dr & v= -\frac{a}{2}e^{-2r/a}
    \end{bmatrix}\\
    = - 2k^2 e^{-2k}  - \frac{4}{a^2} \cdot \frac{a}{2} re^{-2r/a}\at_0^{ka} +\frac{4}{a^2} \cdot \frac{a}{2} \int_0^{ka} e^{-2r/a} \, dr\\
    =- 2k^2 e^{-2k}  - 2k e^{-2k} - \frac{2}{a} \cdot \frac{a}{2} e^{-2r/a} \at_0^{ka}\\
    =\boxed{1-e^{-2k}(2k^2+2k+1)}
\end{gather*}

\subproblem{c}
\begin{gather*}
    P(1)=1-5e^{-2}\\
    P(2)=1-13e^{-4}\\
    P(3)=1-25e^{-6}\\
    \lim_{k \to \infty} P(k) =1
\end{gather*}

\subproblem{d}
\begin{gather*}
    p=\frac{4}{a^3}r^2 e^{-2r/a}\\
    p'=\frac{4}{a^3}(2re^{-2r/a}-\frac{2}{a}r^2e^{-2r/a})\\
    p'=0 \implies r(1-\frac{r}{a})=0 \implies r=\boxed{a}
\end{gather*}
The most likely distance is not $r=0$, as the volume at that value is simply negligible. 

\problem{4}

\subproblem{a}
\begin{gather*}
    x^2+(y-r)^2=r^2 \implies x^2=2ry-y^2\\
    V=\pi \int_0^h 2ry-y^2\, dy\\
    = \pi (ry^2-\frac{y^3}{3})\at_0^h\\
    = \boxed{\pi h^2 (r-\frac{h}{3})}
\end{gather*}

\subproblem{b}
\begin{gather*}
    V=\pi \cdot 10^2\cdot \frac{1}{6}\cdot 10= \frac{500\pi}{3}\\
    \pi h^2(10-\frac{h}{3}) =100 \implies h= 2.797\text{ cm}\\
    \pi h^2(10-\frac{h}{3}) =250 \implies h= 4.850\text{ cm}\\
    \pi h^2(10-\frac{h}{3}) =400 \implies h= 6.843\text{ cm}\\
\end{gather*}

\subproblem{c}
\begin{gather*}
    V=\pi rh^2  - \frac{\pi}{3}h^3\\
    V'= 2\pi r h -\pi h^2\\
    V'' = 2\pi r -2\pi h\\
    V''=0 \implies \boxed{h=r}
\end{gather*}

\problem{5}

\subproblem{a}
\begin{gather*}
    x=6 \implies y=3\\
    y=\frac{x^2}{12} \implies y'= \frac{1}{6}x\\
    V=2\pi \int_0^6  x \sqrt{1+\frac{1}{36}x^2} \, dx\\
    = \begin{bmatrix}
        u=\frac{1}{36}x^2\\
        du= \frac{1}{18} x\, dx
    \end{bmatrix} 36\pi \int_0^1 \sqrt{1+u} \, du\\
    = 36\pi \cdot \frac{2}{3} \cdot (1+u)^\frac{3}{2}\at_0^1\\
    =\boxed{24\pi  (2\sqrt{2} - 1)}
\end{gather*}

\subproblem{b}
\begin{gather*}
    V=\pi \int_0^3 12y \, dy\\
    =6\pi y^2\at_0^3\\
    = \boxed{54\pi}\\
    V=6^2\cdot \pi \cdot 3=108\pi 
\end{gather*}

\subproblem{c}
\begin{gather*}
    x^2+(y-r)^2=r^2 \implies x^2=2ry-y^2\\
    V= l \int _0^h  2\sqrt{2ry-y^2} \, dy\\
    =l \int _0^h  2\sqrt{r^2-(y-r)^2} \, dy\\
    = l \int_0^{\arccos \frac{r-h}{r}}  2r\cos \theta \cdot r\sin \theta\,  d\theta \\
    = l r^2 \int_0^{\arccos \frac{r-h}{r}}  2\sin ^2\theta\,  d\theta \\
    = l r^2 \int_0^{\arccos \frac{r-h}{r}}  1- \cos 2\theta \,  d\theta \\
    = lr^2 (\theta - \frac{\sin 2\theta}{2})\at_0^{\arccos \frac{r-h}{r}}\\
    = lr^2\arccos \paren{\frac{r-h}{r}} - lr^2 \cdot \frac{r-h}{r} \cdot \frac{\sqrt{2rh-h^2}}{r}\\
    = \boxed{l\paren{r^2\arccos \paren{\frac{r-h}{r}} - (r-h)\sqrt{2rh-h^2}}}
\end{gather*}

\subproblem{d}
\begin{gather*}
    V=\pi \cdot 1\cdot 4=4000 \pi\, L\\
    V(0.5)=19.55\%\\
    V=\pi  \implies h=0.596
\end{gather*}
The dipstick is wider near the middle and thinner near the top and the bottom.

\problem{6}

\subproblem{a}
\begin{gather*}
    I= \int_0^R r^2 \cdot \rho \cdot 2\pi rh\, dr\\
    = 2\pi \rho h \cdot \frac{r^4}{4}\at_0^R\\
    =\frac{1}{2}(\pi \rho hR^2)R^2\\
    = \boxed{\frac{1}{2}MR^2}
\end{gather*}

\subproblem{b}
\begin{gather*}
    dm=\rho \cdot 2\pi x \cdot 2\sqrt{R^2-x^2} \, dx\\
    I=  4\pi \rho \int_0^R x^3\sqrt{R^2-x^2} \, dx\\
    = \begin{bmatrix}
        u=R^2-x^2\\
        du=-2x\, dx
    \end{bmatrix} 2\pi \rho \int_{0}^{R^2} (R^2-u)\sqrt{u} \, du\\
    = 2\pi \rho (R^2 \cdot \frac{2}{3}u^\frac{3}{2} - \frac{2}{5}u^\frac{5}{2})\at_0^{R^2}\\
    = \frac{8}{15} \pi \rho R^5\\
    = \boxed{\frac{2}{5}MR^2}
\end{gather*}

\subproblem{c}
\begin{gather*}
    dm= 2\pi R\sigma \, dy\\
    I= \int_{-R}^R (R^2-y^2)2\pi R\sigma \, dy\\
    =4\pi R\sigma \int_0^R R^2-y^2\, dy\\
    =4\pi R\sigma (R^2 y - \frac{y^3}{3})\at_0^R\\
    =\frac{8}{3}\pi R^4 \sigma \\
    =\boxed{\frac{2}{3}MR^2}
\end{gather*}

\subproblem{d}
\begin{gather*}
    Solid \ ball > Cylinder > Hollow \ ball
\end{gather*}
The answer solely depends on the moment of inertia, not the mass.

\problem{7}

\subproblem{a}
\begin{gather*}
    x=3\implies y = 2\sqrt{3}\\
    y=\frac{2}{3}x^\frac{3}{2} \implies y'=\sqrt{x}\\
    L= \int_0^3 \sqrt{1+x} \, dx\\
    =   \frac{2}{3}(1+x)^\frac{3}{2}\at_0^3\\
    = \boxed{\frac{14}{3}}
\end{gather*}

\subproblem{b}
\begin{gather*}
    S= 2\pi \int_0^3 x \sqrt{1+x} \, dx\\
    = 2\pi \int_1^4(u-1)\sqrt{u} \, du\\
    = 2\pi \int_1^4 u^\frac{3}{2}-\sqrt{u} \, du\\
    = 2\pi (\frac{2}{5}u^\frac{5}{2} - \frac{2}{3}u^\frac{3}{2})\at_1^4\\
    = \boxed{\frac{232\pi}{15}}
\end{gather*}

\subproblem{c}
\begin{gather*}
    V=2\pi \int_0^3  x (2\sqrt{3}-\frac{2}{3}x^\frac{3}{2})\, dx\\
    = 2\pi \int_0^3 2\sqrt{3}x - \frac{2}{3}x^\frac{5}{2} \, dx\\
    = 2\pi (\sqrt{3}x^2 - \frac{2}{3}\cdot \frac{2}{7} x^\frac{7}{2})\at_0^3\\
    =2\pi (9\sqrt{3} - \frac{36\sqrt{3}}{7} )\\
    =\boxed{\frac{54\sqrt{3}}{7}\pi }
\end{gather*}

\subproblem{d}
\begin{gather*}
    \frac{3}{2}y=x^\frac{3}{2} \implies x^2= (\frac{3}{2}y)^\frac{4}{3}\\
    V=\pi \int_0^{2\sqrt{3}}(\frac{3}{2}y)^\frac{4}{3}\, dy\\
    = \pi \cdot \frac{3}{2}^\frac{4}{3} \cdot \frac{3}{7}y^\frac{7}{3} \at_0^{2\sqrt{3}}\\
    = \boxed{\frac{54\sqrt{3}}{7}\pi }
\end{gather*}

\subproblem{e}
\begin{gather*}
    V=\pi \cdot \frac{3}{2}^\frac{4}{3} \cdot \frac{3}{7}h^\frac{7}{3}\\
    V(\frac{1}{2}h)=\frac{1}{2^{7/3}}V
\end{gather*}
The area is much less than 1/2 as the bowl is more narrow near the bottom.

\problem{8}

\subproblem{c}
\begin{gather*}
    h=\frac{2d}{\sqrt{3}} \implies d = \frac{\sqrt{3}h}{2}\\
    V_{max} = \pi (\frac{d^2}{4}h-\frac{1}{16}h^3)=\frac{\pi h^3}{8}\\
    V=\frac{\pi}{16}(3h^2 \cdot h(1+\varepsilon)- h^3(1+\varepsilon)^3)\\
    = \frac{\pi h^3}{8}(1-\frac{3}{2}\varepsilon^2 -\frac{1}{2}\varepsilon^3)\\
    = \boxed{ V_{max}(1-\frac{3}{2}\varepsilon^2 -\frac{1}{2}\varepsilon^3)}
\end{gather*}





















\end{document}

